\documentclass[17pt,a4paper,landscape]{extarticle}
\usepackage[margin=2.35cm,top=0.12cm,left=1.05cm]{geometry}
\usepackage{amsmath}
\usepackage{amssymb}
\usepackage{tasks}
\usepackage{tikz}
\usepackage{xcolor}
\usepackage[sfdefault,lf]{FiraSans}
\renewcommand{\familydefault}{\sfdefault}
\setlength{\parindent}{0pt}
\pagestyle{empty}
\settasks{label=\textbf{\arabic*.}, label-width=2.2em, item-indent=2.95em, column-sep=2.1cm, after-item-skip=4.7em}
\renewcommand{\arraystretch}{1.15}
\everymath{\displaystyle}

\begin{document}
\sffamily
\boldmath

\boldmath

{\LARGE \textbf{Understanding Cartesian plane and co-ordinates — Excellence}}
\begin{tasks}(3)
\task {Write the co-ordinates of point \mbox{$A$} if it is 4 left and 5 up. Then describe its location using a quadrant.}
\task {Write the co-ordinates of point \mbox{$B$} if it is on the \mbox{$y$}-axis and 4 down.}
\task {Write the co-ordinates of point \mbox{$C$} if it is 3 right and 2 down. Then describe its location using a quadrant.}
\task {Write the co-ordinates of point \mbox{$D$} if it is on the \mbox{$x$}-axis and 1 left.}
\task {Plot \mbox{$A(-4,3)$}, \mbox{$B(4,3)$}, \mbox{$C(4,-3)$} and \mbox{$D(-4,-3)$}. Join them in order. What shape is formed?}
\task {Point \mbox{$P$} is at \mbox{$(2,-1)$}. Point \mbox{$Q$} has the same \mbox{$x$}-value but the opposite \mbox{$y$}-value. Write the co-ordinates of \mbox{$Q$}.}
\task {Point \mbox{$M$} is at \mbox{$(-3,4)$}. Point \mbox{$N$} has the opposite \mbox{$x$}-value but the same \mbox{$y$}-value. Write the co-ordinates of \mbox{$N$}.}
\task {Two points lie on the horizontal line through \mbox{$(0,2)$}. One is 5 units left of the \mbox{$y$}-axis and one is 4 units right of the \mbox{$y$}-axis. Write both co-ordinates.}
\task {Two points lie on the vertical line through \mbox{$(-1,0)$}. One is 3 units above the \mbox{$x$}-axis and one is 4 units below the \mbox{$x$}-axis. Write both co-ordinates.}
\task {Which pair has the same distance from the origin along an axis: \mbox{$A(-5,0)$}, \mbox{$B(0,5)$}, \mbox{$C(4,0)$}, \mbox{$D(0,-4)$}?}
\task {Explain why “Y to the Sky” helps students place the point \mbox{$(-2,5)$} correctly.}
\task {One point is 3 units right and 2 units down from the origin. Another is 3 units left and 2 units up from the origin. Write both co-ordinates and name their quadrants.}
\task {Plot \mbox{$E(-2,4)$}, \mbox{$F(2,4)$}, \mbox{$G(2,-1)$} and \mbox{$H(-2,-1)$}. Which side lengths can you read directly from the grid?}
\task {Make up one point in Quadrant II and one point in Quadrant IV that have the same distance from the \mbox{$x$}-axis.}
\task {Make up three different points whose \mbox{$x$}-values add to \mbox{$0$}.}
\task {Explain the difference between \mbox{$(4,0)$} and \mbox{$(0,4)$}.}
\task {A student writes \mbox{$(-3,2)$} as \mbox{$(2,-3)$}. Describe exactly how the location changes on the grid.}
\task {Write one point that is on an axis but not at the origin, and explain how you know it is on an axis.}
\task {Plot any four points that make a rectangle with sides parallel to the axes. Write the co-ordinates of your four points.}
\end{tasks}

\end{document}
